% ---------
% --------
%!TEX TS-program = pdflatex
%!TEX encoding = UTF-8 Unicode
%!TEX spellcheck = English (Aspell)

\documentclass[11pt]{article}
\usepackage{jeffe,handout,graphicx}
\usepackage{fourier-orns}
\usepackage{mathtools}
\usepackage[normalem]{ulem} % either use this (simple) or
\usepackage{tikz}
\pagestyle{empty}

\allowdisplaybreaks

\begin{document}


\begin{center}
\Large\textbf{CSCI 332, Fall 2025}%
\\
\LARGE\textbf{In-class Activity}%
\end{center}
For each of the three problems, answer the questions at the bottom of the sheet.
\begin{enumerate}
   \item 
You are climbing a staircase. It takes n steps to reach the top. Each time you
can either climb 1 or 2 steps. In how many distinct ways can you climb to the
top?
\item
You are given an integer array cost where $cost[i]$ is the cost of ith step on a
staircase. Once you pay the cost, you can either climb one or two steps. You
can either start from the step with index 0, or the step with index 1. Return the
minimum cost to reach the top of the staircase.
\item
In your country, there are only coins of denominations 1, 3, and 4. Given a
nonnegative integer $n$, what is the fewest number of coins you can use to
represent $n$?
\end{enumerate}

\subsection*{Questions to answer for each problem: }
\begin{enumerate}
   \item Write the definition of the recursive subproblems you will solve. Make sure to 
   give these a descriptive name.
   \item Write the recurrence to solve your recursive subproblems.
   \item In what direction will you fill your memoization array?
   \item Write the pseudocode for the memoized algorithm.
\end{enumerate}

%%%%%%%%%%%%%%%%%%%%
\end{document}
